Para descrever a dinâmica de atitude completa de uma espaçonave, deve-se integrar todas as equações cinemáticas e dinâmicas.

\subsection{Simulações numéricas}

Enquanto que as equações de Euler
\eqref{eq:euler_corpo_rigido}
fornecem a evolução da velocidade angular em função dos torques aplicados, as equações cinemáticas
\eqref{eq:definicao_qdot} propagam a orientação da espaçonave através dos quaternions unitários. 
% A atitude relativa \eqref{eq:quat_relativo} define a orientação da espaçonave visitante em relação a espaçonave alvo, portanto, transformando-as na base para o cálculo dos erros de orientação empregados nas leis de controle.

O processo de simulação numérica consiste em resolver as equações apresentadas, de forma acoplada:

\begin{equation}
\begin{cases}
\mathbf{I}(t)\,\dot{\vec{\omega}}+\vec{\omega}\times(\mathbf{I}(t)\,\vec{\omega})
= \vec\tau, \\[6pt]
\dot{\vec{q}}
= \tfrac{1}{2}\,\Omega(\vec{q})\,\vec{\omega}, \\[6pt]
q_{rel}
= q_t^{-1} \otimes q_c,
\end{cases}
\label{eq:sistema_acoplado}
\end{equation}
onde $\mathbf{I}(t)$ é a matriz de inércia equivalente, $\vec{\omega}$ é a velocidade angular, $\vec\tau$ os torques aplicados, e $\Omega(\vec{q})$, sendo o operador cinemático.

% A integração numérica é realizada em tempo contínuo, utilizando métodos de passo variável, como o \texttt{ode45}, que é adequado para sistemas dinâmicos não lineares. 
% Na simulação, os quaternions são constantemente normalizados para mitigar erros numéricos de arredondamento e manter a unidade da norma. 
% Sendo possível avaliar o desempenho do controlador em diferentes condições iniciais e durante as manobras acopladas do manipulador robótico.